\documentclass[10pt,a4paper]{amsart}
\usepackage[margin=19mm]{geometry}
\usepackage{amsmath,amssymb,mathtools}
\usepackage[scr=rsfs,cal=euler]{mathalfa}
\usepackage{xfrac}
\usepackage[unicode,hidelinks,bookmarks=false]{hyperref}
\newcommand{\ie}{i.e.\ }
\newcommand{\cf}{cf.\ }
\newcommand{\resp}{respectively\ }
\newcommand{\Subsection}{\subsection}
\providecommand{\nobreakdash}{\nobreak\hbox{-}}
\setlength{\emergencystretch}{2em}
\multlinegap0pt
\usepackage{bm}
\usepackage{amssymb}
\usepackage{algorithm}\usepackage{algpseudocode}
\usepackage[shortlabels]{enumitem}
\usepackage{booktabs}
\newtheorem{proposition}{Proposition}
\def\eqref#1{equation~\ref{#1}}
\title{From Adam to Adam-Like Lagrangians: Second-Order Nonlocal Dynamics}
\author{Carlos Heredia}
\date{Source: arXiv:2602.09101v1 (CC BY 4.0)}
\begin{document}
\maketitle

\noindent\emph{Source and licence.} From Adam to Adam-Like Lagrangians: Second-Order Nonlocal Dynamics. Version arXiv:2602.09101v1 (CC BY 4.0), distributed by arXiv under the Creative Commons Attribution 4.0 International licence, http://creativecommons.org/licenses/by/4.0/, as recorded in the arXiv licence metadata for this exact version and shown on the abstract page https://arxiv.org/abs/2602.09101v1. Full licence notice: \texttt{licenses/arxiv-2602.09101-CC-BY-4.0.txt}.\par\smallskip

\noindent\textit{Editorial scope.} Counted unit: the article's proposition prop:alpha-recovery (source0.tex lines 919--972). The article's complete original proof of it is delivered with it (source0.tex lines 974--1131, one \texttt{proof} environment of 158 lines). No mathematics is written, repaired or re-derived: the statement and the proof are the article's own lines, in the article's own notation, and no retained line is substituted or re-flowed.  Retained uncounted as the source-local context the proof needs: lines 693--729.  Nothing was omitted from any retained span, so the package carries no omission record.  No retained line contains a citation, so the wrapper carries no bibliography and no bibliography entry.  No retained line contains a figure environment, an image-inclusion command or drawing code, so no image is delivered and the package declares no asset dependency.  Editorial formatting only, in addition: an AMS \texttt{amsart} preamble that restates the article's own declarations and macros used by the retained lines, namely the environments \texttt{proposition} on separate counters, together with the standard packages those lines need.  The retained environments print in this document as Proposition 1 in order of appearance, whereas the article prints the same items with the numbering of its own full document.  The complete native source of the article as distributed in the arXiv e-print for this exact version is bundled unchanged as \texttt{sources/194295/source0.tex}.
\begin{proposition}\label{prop:Adam2}
Let $\alpha>0$ be small and consider the continuous limit $t\approx k\,\alpha$ with the second-order time expansion
\[
\theta^i_{k+1} \to \theta^i(t+\alpha)\;=\;\theta^i(t)+\alpha\,\dot\theta^i(t)+\tfrac{\alpha^2}{2}\,\ddot\theta^i(t)\,+\mathcal{O}(\alpha^3).
\]
With the initial data $m^i(0)=\dot m^i(0)=0$ and $v(0)=\dot v(0)=0$, the Adam update admits the inertial ODE with linear friction
\begin{equation}\label{eq:inertial}
\frac{\alpha}{2}\,\ddot\theta^i(t)+\dot\theta^i(t)\;=\;-\eta(t+\alpha)\,T^i(\theta,t+\alpha),
\qquad
T^i(\theta,t)\;:=\;\frac{m^i(\theta,t)}{\sqrt{v(\theta,t)}+\varepsilon(t)}, 
\end{equation}
where the continuous bias-correction factors are
\[
\eta(t)\;:=\;\frac{\sqrt{\,1-\beta_2^{\,t/\alpha}\,}}{\,1-\beta_1^{\,t/\alpha}\,},\qquad
\varepsilon(t)\;:=\;\epsilon\,\sqrt{\,1-\beta_2^{\,t/\alpha}\,}.
\]
The continuous moments are causal convolutions
\[
m^i(\theta,t)\;=\;\frac{1-\beta_1}{\alpha}\int_{0}^{t} K_{\beta_1}(t-\tau)\;\partial^i f\!\big(\theta(\tau)\big)\,d\tau,\quad
v(\theta,t)\;=\;\frac{1-\beta_2}{\alpha}\int_{0}^{t} K_{\beta_2}(t-\tau)\;\|\nabla f(\theta(\tau))\|^2\,d\tau,
\]
with kernels $K_\beta(s)$ (for $s:=t-\tau\ge 0$) given case-wise by
\begin{itemize}
\item[\textbullet] Critical damping $\beta=\tfrac12$:
\[
K_{1/2}(s)=\frac{2\,s}{\alpha}\,e^{-s/\alpha}.
\]
\item[\textbullet] Overdamped regime $\beta>\tfrac12$. With $\kappa:=\sqrt{\,2\beta-1\,}$,
\[
K_\beta(s)\;=\; \frac{2}{\kappa}\,e^{-s/\alpha}\,\sinh\!\Big(\frac{\kappa}{\alpha}s\Big)\,.
\]
\item[\textbullet] Underdamped regime $\beta<\tfrac12$. With $\rho:=\sqrt{\,1-2\beta\,}$,
\[
K_\beta(s)=\frac{2}{\rho}\,e^{-s/\alpha}\,\sin\!\Big(\frac{\rho}{\alpha}s\Big).
\]
\end{itemize}
\end{proposition}

\begin{proposition}\label{prop:alpha-recovery}
Let $T>0$, $0<\delta<T$, and $\alpha_0\in(0,\delta)$. For $a\in\{1,2\}$ define
\[
H^{(1)}_a(\tau):=\frac{1-\beta_a}{\alpha}\exp\!\Big(-\frac{1-\beta_a}{\alpha}\tau\Big),
\qquad
H^{(2)}_a(\tau):=\frac{1-\beta_a}{\alpha}\,K_{\beta_a}(\tau),
\qquad \tau\ge 0,
\]
where $K_\beta$ is given in Proposition~\ref{prop:Adam2}. Let $g=\nabla f$ and set, for $j\in\{1,2\}$,
\[
m^{(j)}(t):=\int_0^t H^{(j)}_1(\tau)\,g(\theta(t-\tau))\,d\tau,
\qquad
v^{(j)}(t):=\int_0^t H^{(j)}_2(\tau)\,\|g(\theta(t-\tau))\|^2\,d\tau .
\]

Assume $\theta\in C^3([0,T];\mathbb R^d)$ and that the maps $t\mapsto g(\theta(t))$ and $t\mapsto \|g(\theta(t))\|^2$ have bounded second derivatives on $[0,T]$.
Assume moreover that
\[
\sup_{t\in[\delta-\alpha_0,\,T]}\|\ddot\theta(t)\|\le C_T
\]
for some $C_T$ independent of $\alpha$. Then there exist $\varepsilon_0>0$ (depending on $\delta$ but not on $\alpha$) such that $\varepsilon(t)\ge \varepsilon_0$ for all $t\in[\delta,T]$. Furthermore, there exists $C>0$, independent of $\alpha$, such that for all $t\in[\delta,T]$,
\[
\|m^{(2)}(t)-m^{(1)}(t)\|
\le C\Big(\frac{\alpha}{1-\beta_1}\Big)^{2},
\qquad
\bigl|v^{(2)}(t)-v^{(1)}(t)\bigr|
\le C\Big(\frac{\alpha}{1-\beta_2}\Big)^{2}.
\]
Consequently, uniformly for $t\in[\delta,T]$,
\[
\frac{m^{(2)}(t)}{\sqrt{v^{(2)}(t)}+\varepsilon(t)}
=
\frac{m^{(1)}(t)}{\sqrt{v^{(1)}(t)}+\varepsilon(t)}
+\mathcal{O}\!\left(\max\left\{\Big(\tfrac{\alpha}{1-\beta_1}\Big)^2,\ \tfrac{\alpha}{1-\beta_2}\right\}\right).
\]

Finally, if $0<\alpha<\alpha_0$ and $\theta$ satisfies
\[
\frac{\alpha}{2}\ddot\theta(t)+\dot\theta(t)
=
-\eta(t+\alpha)\,
\frac{m^{(2)}(t+\alpha)}{\sqrt{v^{(2)}(t+\alpha)}+\varepsilon(t+\alpha)},
\]
then, uniformly for $t\in[\delta-\alpha,\,T-\alpha]$,
\[
\dot\theta(t)
=
-\eta(t+\alpha)\,
\frac{m^{(1)}(t+\alpha)}{\sqrt{v^{(1)}(t+\alpha)}+\varepsilon(t+\alpha)}
+\mathcal{O}(\rho),
\qquad
\rho:=\max\left\{\alpha,\Big(\tfrac{\alpha}{1-\beta_1}\Big)^2,\ \tfrac{\alpha}{1-\beta_2}\right\}.
\]
\end{proposition}

\begin{proof}
We work with the truncated causal convolution
\[
(H*\varphi)(t):=\int_0^t H(\tau)\,\varphi(t-\tau)\,d\tau.
\]
For $a\in\{1,2\}$ and $j\in\{1,2\}$, define the full moments
\[
M^{(j)}_{a,k}:=\int_0^\infty \tau^k\,H^{(j)}_a(\tau)\,d\tau,\qquad k=0,1,2,
\]
and the truncated moments
\[
M^{(j)}_{a,k}(t):=\int_0^t \tau^k\,H^{(j)}_a(\tau)\,d\tau,\qquad t\in[0,T].
\]

For $H^{(1)}_a(\tau)=\lambda_a e^{-\lambda_a\tau}$ with $\lambda_a=(1-\beta_a)/\alpha$,
\[
M^{(1)}_{a,0}=1,\qquad
M^{(1)}_{a,1}=\lambda_a^{-1},\qquad
M^{(1)}_{a,2}=2\,\lambda_a^{-2}.
\]
For $H^{(2)}_a(\tau)=\lambda_a\,K_{\beta_a}(\tau)$, one obtains using the closed forms in Proposition~\ref{prop:Adam2}:
\[
M^{(2)}_{a,0}=1,\qquad
M^{(2)}_{a,1}=\lambda_a^{-1},\qquad
M^{(2)}_{a,2}=(1+\beta_a)\,\lambda_a^{-2}.
\]

Since the kernels decay exponentially on the scale $\alpha$, there exist constants $c,C>0$
(independent of $\alpha$) such that for $k=0,1$ and all $t\in[\delta,T]$,
\[
\bigl|M^{(j)}_{a,k}(t)-M^{(j)}_{a,k}\bigr|
=\int_t^\infty \tau^k |H^{(j)}_a(\tau)|\,d\tau
\le C\,e^{-c\,t/\alpha}
\le C\,e^{-c\,\delta/\alpha}.
\]
In particular, since $M^{(1)}_{a,0}=M^{(2)}_{a,0}=1$ and $M^{(1)}_{a,1}=M^{(2)}_{a,1}=\lambda_a^{-1}$,
we have
\[
\bigl|M^{(1)}_{a,0}(t)-M^{(2)}_{a,0}(t)\bigr|
+\bigl|M^{(1)}_{a,1}(t)-M^{(2)}_{a,1}(t)\bigr|
\le C\,e^{-c\,\delta/\alpha},
\qquad t\in[\delta,T].
\]

Let $H^{(1)},H^{(2)}$ be causal kernels and let $\varphi\in C^3([0,T])$.
For each $t\in[0,T]$ and $0\le \tau\le t$, Taylor's theorem gives
\[
\varphi(t-\tau)=\varphi(t)-\tau \dot\varphi(t)+\frac{\tau^2}{2}\,\ddot \varphi(t) + \mathcal{O}(\tau^3).
\]
Hence, 
\begin{align*}
(H^{(1)}*\varphi)(t)-(H^{(2)}*\varphi)(t)
&=\varphi(t)\bigl(M^{(1)}_0(t)-M^{(2)}_0(t)\bigr)
-\dot\varphi(t)\bigl(M^{(1)}_1(t)-M^{(2)}_1(t)\bigr)\\
&\quad+\frac12 \ddot\varphi(t)\int_0^t \tau^2\Bigl(H^{(1)}(\tau)-H^{(2)}(\tau)\Bigr)\,d\tau + \mathcal{O}(\tau^3).
\end{align*}
Taking absolute values and using $\|\ddot\varphi\|_{L^\infty([0,T])}<\infty$ yields
\[
\bigl|(H^{(1)}*\varphi)(t)-(H^{(2)}*\varphi)(t)\bigr|
\le C\,e^{-c\,\delta/\alpha}
+\frac{\|\ddot\varphi\|_{L^\infty([\delta,T])}}{2}
\left(\int_0^\infty \tau^2|H^{(1)}(\tau)|\,d\tau+\int_0^\infty \tau^2|H^{(2)}(\tau)|\,d\tau\right),
\]
uniformly for $t\in[\delta,T]$. For the kernels $H^{(j)}_a$, the second-moment integrals scale like $\lambda_a^{-2}$, hence
\[
\bigl|(H^{(1)}*\varphi)(t)-(H^{(2)}*\varphi)(t)\bigr|
\le C\left(\frac{\alpha}{1-\beta_a}\right)^2 + C\,e^{-c\,\delta/\alpha},
\qquad t\in[\delta,T].
\]
For $\alpha$ sufficiently small (depending on $\delta$), the exponential term is dominated by the
polynomial term, so the bound simplifies to $C(\alpha/(1-\beta_a))^2$.

Now, we apply the previous estimate componentwise with $\varphi_i(t)=g_i(\theta(t))$ and with $\varphi(t)=\|g(\theta(t))\|^2$. By the assumptions, $\varphi_i''$ and $\varphi''$ are bounded on $[0,T]$. Using the bounds on $M^{(j)}_{a,2}$ we deduce
\[
\|m^{(2)}(t)-m^{(1)}(t)\|
=
\mathcal O\!\Big(\Big(\tfrac{\alpha}{1-\beta_1}\Big)^{\!2}\Big),
\qquad
|v^{(2)}(t)-v^{(1)}(t)|
=
\mathcal O\!\Big(\Big(\tfrac{\alpha}{1-\beta_2}\Big)^{\!2}\Big),
\]
uniformly on $[\delta,T]$.

Let $F(x,y)=x/(\sqrt{y}+\varepsilon)$. For $y\ge 0$, $y+\delta y \geq 0$ and $\varepsilon\ge\varepsilon_0$, we can deduce that:
\[
|F(x+\delta x,y+\delta y)-F(x,y)|
\le
\frac{|\delta x|}{\varepsilon_0}
+
\,\frac{|x|}{\varepsilon_0^2}\,\sqrt{|\delta y|}\,.
\]
Indeed,
\[
\frac{x+\delta x}{\sqrt{y+\delta y}+\varepsilon}-\frac{x}{\sqrt{y}+\varepsilon}
=\frac{\delta x}{\sqrt{y+\delta y}+\varepsilon} + x\left[\frac{1}{\sqrt{y+\delta y} + \varepsilon} - \frac{1}{\sqrt{y}+\varepsilon} \right]
\]
The first term is bounded by $|\delta x|/\varepsilon_0$. For the second, we use the fact that
\[
|\sqrt{y+\delta y}-\sqrt{y}| \leq \frac{|\delta y|}{\sqrt{y+\delta y} + \sqrt{y}}\leq \sqrt{|\delta y|}\,,
\]
therefore, it becomes bounded by
\[
\left|x\!\left(\frac{1}{\sqrt{y+\delta y}+\varepsilon}-\frac{1}{\sqrt{y}+\varepsilon}\right)\right| \leq \frac{|x|}{\varepsilon^2_0} \sqrt{|\delta y|}\,.
\]

Thus, taking $x=m^{(1)}(t)$, $y=v^{(1)}(t)$, $\delta x=m^{(2)}(t)-m^{(1)}(t)$ and $\delta y=v^{(2)}(t)-v^{(1)}(t)$, and assuming $v^{(j)}(t)\ge0$, we have\footnote{Note that $y+\delta y=v^{(2)}(t)\ge 0$ by assumption.}
\[
\delta x=\mathcal{O}\!\Big(\Big(\tfrac{\alpha}{1-\beta_1}\Big)^{\!2}\Big),\qquad
\delta y=\mathcal{O}\!\Big(\Big(\tfrac{\alpha}{1-\beta_2}\Big)^{\!2}\Big).
\]
Notice that, since $[0,T]$ is compact and $g\circ\theta$ is continuous, there exists
\[
G_T:=\sup_{s\in[0,T]}\|g(\theta(s))\|<\infty.
\]
Therefore, using $H^{(1)}_1\ge 0$ and $\int_0^t H^{(1)}_1(\tau)\,d\tau\ \leq 1$, we obtain
\[
\|m^{(1)}(t)\|
\le \int_0^t H^{(1)}_1(\tau)\,\|g(\theta(t-\tau))\|\,d\tau
\le G_T\int_0^t H^{(1)}_1(\tau)\,d\tau
\leq G_T,
\]
so that $|x|=\|m^{(1)}(t)\|$ is uniformly bounded and can be absorbed into the implicit constant. Hence,
\begin{equation}\label{eq:diff_moments_order}
  \frac{m^{(2)}(t)}{\sqrt{\,v^{(2)}(t)}+\varepsilon(t)\,}
-\frac{m^{(1)}(t)}{\sqrt{\,v^{(1)}(t)}+\varepsilon(t)\,}
=
\mathcal{O}\!\left(\max\left\{\Big(\tfrac{\alpha}{1-\beta_1}\Big)^2,\ \tfrac{\alpha}{1-\beta_2}\right\}\right), \qquad t\in[\delta, T].
\end{equation}


Finally, we rewrite the accelerated shifted equation as
\[
\dot\theta(t)
=
-\eta(t+\alpha)\,
\frac{m^{(2)}(t+\alpha)}{\sqrt{v^{(2)}(t+\alpha)}+\varepsilon(t+\alpha)}
-\frac{\alpha}{2}\ddot\theta(t).
\]
By assumption $\sup_{t\in[\delta-\alpha_0,\,T]}\|\ddot\theta(t)\|\le C_T$, hence
\[
\Big\|\tfrac{\alpha}{2}\ddot\theta(t)\Big\|\le \tfrac{\alpha}{2}C_T=\mathcal O(\alpha)
\quad\text{uniformly for }t\in[\delta-\alpha,\,T-\alpha],
\]
for all $0<\alpha<\alpha_0$ with $\alpha_0\in(0,\delta)$. Together with the kernel comparison---\eqref{eq:diff_moments_order}---, this yields, upon substitution into the accelerated dynamics,
\[
\dot\theta(t)=
-\eta(t+\alpha)\,
\frac{m^{(1)}(t+\alpha)}{\sqrt{v^{(1)}(t+\alpha)}+\varepsilon(t+\alpha)}
+\mathcal O(\rho),
\quad t\in[\delta-\alpha,T-\alpha]\,,
\]
where
\[
\rho:=\max\left\{\alpha,\Big(\tfrac{\alpha}{1-\beta_1}\Big)^2,\ \tfrac{\alpha}{1-\beta_2}\right\}\,.
\]
In particular, if $\beta_1,\beta_2$ are fixed and $\alpha\to0$, then $\left(\alpha/(1-\beta_1)\right)^2 = \mathcal O(\alpha^2)\subset \mathcal{O}(\alpha)$ and $\alpha/(1-\beta_2) = \mathcal O(\alpha)$, so the remainder can be written as $\mathcal O(\alpha)$.
\end{proof}

\end{document}
